\section*{Замечания рецензентов и ответы авторов}
\noindent Заявка 90. DAMDID 2026.

\paragraph{Рецензент 1: пояснения и примеры.}
\emph{Замечание, краткий пересказ.} Статья содержит необходимые исследовательские разделы, однако смысл углублённой доразметки неясен. Рецензенту трудно следить за английским изложением; для объяснения основных понятий предложено добавить конкретные примеры.

\emph{Ответ.} В аннотацию и введение добавлено определение углублённой доразметки как расширения оценок на несколько ответов-кандидатов в выбранном пуле. Охват разметки отделён от глубины нейросети и исчерпывающей разметки корпуса. Во введении также поясняется, что более полные оценки меняют измерение качества фиксированного ранжирования, а не улучшают саму поисковую модель. В разделе~\ref{sec:example} и таблице~\ref{tab:annotation-example} приведены реальный запрос MS MARCO, идентификаторы фрагментов, сокращённое содержание и метки до и после доразметки. Пример показывает, почему правильный, но неразмеченный ответ не получает положительного вклада при неполных qrels. Пояснения включены в обе языковые версии.

\paragraph{Рецензент 2: общая оценка.}
\emph{Замечание, краткий пересказ.} Рецензент отмечает актуальность улучшения наборов данных, ценность результатов и обоснованность подхода и поддерживает публикацию полной статьи.

\emph{Ответ.} Благодарим за положительную оценку. Конкретных дополнительных изменений рецензент не запросил.

\paragraph{Рецензент 3: другие провайдеры и универсальность.}
\emph{Замечание, краткий пересказ.} Есть ли аналоги MS MARCO у Google, Yandex и Yahoo? Насколько универсальна MS MARCO с учётом различий в использовании поисковых систем между странами и социальными группами? Предложено обосновать выбор коллекции во введении.

\emph{Ответ.} Во введение добавлены два абзаца о Natural Questions от Google, Learning to Rank Challenge от Yahoo! и Personalized Web Search Challenge от Yandex со ссылками на публикации или страницы наборов данных. Разграничены ответы на вопросы по текстам, векторы признаков для ранжирования и анонимизированные журналы взаимодействий; эти ресурсы не представлены как взаимозаменяемые наборы для ранжирования фрагментов. Выбор MS MARCO объяснён доступностью текстов, неполной разметкой, устоявшимися средствами оценки и предобученными моделями. Прямо указано, что различия между странами и социальными группами не измерялись, а коллекция не является универсальной выборкой информационных потребностей.
